\chapter{Implementación de las simulaciones}\label{ap:simulacion}

Todas las simulaciones del capítulo~\ref{chap:desempeno} siguen el mismo esquema de tiempo discreto, que reproduce un controlador digital con retenedor de orden cero. En cada periodo de muestreo $T=1$~ms se ejecutan los pasos siguientes.

\begin{enumerate}
  \item \textbf{Planta.} Se integra el modelo no lineal durante un periodo con el voltaje $u$ constante, con \texttt{ode45} de MATLAB. Con atascamiento-deslizamiento, la aceleración se calcula con el modelo de Karnopp: si $|v|>DV$, el imán desliza con fricción viscosa $m c_1 v$; si $|v|\le DV$ y la fuerza neta $F_e=u/[b(a-x)^4]-mg$ cumple $|F_e|\le F_s$, permanece atascado; en otro caso se despega con aceleración $(F_e - F_s\,\operatorname{sign}F_e)/m$.
  \item \textbf{Sujeción de atascamiento.} Al final del periodo, si $|v|<DV$ y $|F_e|\le F_s$, la velocidad se fija en cero. Esto evita que el integrador numérico deje una velocidad residual dentro de la banda de Karnopp.
  \item \textbf{Controlador.} Se calcula el error $e=r-x$ y se integra el modelo en espacio de estados del controlador durante el mismo periodo con $e$ constante.
  \item \textbf{Saturación.} La salida del controlador se acota al intervalo $[0;\,3{,}5]$~V.
  \item \textbf{Zona muerta.} Si el imán está en movimiento y el nuevo voltaje difiere del último aplicado en menos de $+0{,}025$~V o $-0{,}02$~V, se conserva el último voltaje aplicado.
\end{enumerate}

El controlador se inicializa en su estado estacionario: se resuelve $(A_d - I)\,x_c = 0$ junto con $C x_c = u_e(x_0)$, donde $A_d$ es la matriz del controlador discretizado, de modo que en $t=0$ entrega exactamente el voltaje de equilibrio de la posición inicial y el levitante parte del reposo. Para el PD con prealimentación de la sección~\ref{sec:analisis_ciclo} se suma al voltaje del controlador el voltaje de equilibrio $u_e(r)$ de la referencia.

Los barridos paramétricos (rango viable del actuador y familias de controladores) se ejecutaron con una implementación equivalente en la GPU, con el controlador discretizado de forma exacta y la planta integrada por Runge-Kutta de cuarto orden con diez subpasos por periodo. Sus resultados se validaron contra la implementación con \texttt{ode45} en casos representativos, con diferencias menores al 8\,\%.

La tabla~\ref{tab:scripts} relaciona los programas con las secciones de la tesis.

\begin{table}[ht]
\centering
\caption{Programas que generan los resultados de la tesis.}
\label{tab:scripts}
\small
\begin{tabular}{|l|p{0.48\textwidth}|}
\hline
Programa & Resultados \\ \hline
\texttt{evidencia\_P8.m} & Lazo abierto, modelo lineal con AD, configuración estable \\ \hline
\texttt{diseno\_PII\_tikz.m} & Diagramas de Bode, Nyquist y respuesta al escalón \\ \hline
\texttt{regulacion\_sinAD\_P3.m} & Modelo no lineal sin atascamiento-deslizamiento \\ \hline
\texttt{regulacion\_P4.m} & Regulación \\ \hline
\texttt{seguimiento\_P5.m} & Seguimiento senoidal y trapezoidal \\ \hline
\texttt{comparar\_PI\_PII\_P6.m} & Comparación PI frente a PII \\ \hline
\texttt{familias\_R15.m}, \texttt{simular\_R15.py} & Familias de controladores PI y PII \\ \hline
\texttt{analisis\_ciclo.m}, \texttt{exportar\_ciclo.py} & Análisis del ciclo límite \\ \hline
\texttt{barrido\_3v5.py} & Rango viable con $u\le3{,}5$~V \\ \hline
\end{tabular}
\end{table}
